% Shared colors and styles of the figures. Every figure reads its data from figures/data/*.csv, written by
% experiments/make_figures.py (comma separated, header row, nan for missing values: pgfplots skips those points).
\definecolor{cA}{HTML}{2a78d6}\definecolor{cB}{HTML}{eb6834}\definecolor{cC}{HTML}{1baf7a}%
\definecolor{cD}{HTML}{4a3aa7}\definecolor{cE}{HTML}{c03b7a}\definecolor{cF}{HTML}{8a6d1f}\definecolor{cGray}{HTML}{8a8984}%
\definecolor{lvA}{HTML}{c6dbf5}\definecolor{lvB}{HTML}{93bbec}\definecolor{lvC}{HTML}{5e98e0}%
\definecolor{lvD}{HTML}{2a78d6}\definecolor{lvE}{HTML}{174a8c}\definecolor{lvF}{HTML}{1baf7a}%
%
% \namedbars{csv file}{ylabel}{title}: lower bounds of the named states at every level, divided by the reference `ref`
% (printed above each state at height `label_y`), with the two upper bounds as markers. Columns: see named_csv in
% make_figures.py. The state positions (`pos`) leave a gap at x = 10 between the PPT and the NPT states.
\providecommand{\namedbar}[3]{% fill, column, offset
  \addplot[ybar, bar width=0.13, bar shift=0pt, draw=none, fill=#1, area legend]
    table[x expr=\thisrow{pos}+#3, y expr=\thisrow{#2}/\thisrow{ref}] {\namedfile};}
\providecommand{\namedbars}[3]{%
  \def\namedfile{#1}%
  \begin{tikzpicture}
    \begin{axis}[width=\textwidth, height=0.29\textwidth, ymin=0, ymax=1.38, xmin=0.3, xmax=14.7,
      ylabel={#2}, ytick={0,0.25,0.5,0.75,1}, xtick=data, xticklabels from table={#1}{state},
      table/col sep=comma, xticklabel style={rotate=35, anchor=north east, font=\scriptsize, inner sep=1pt},
      tick label style={font=\footnotesize}, label style={font=\small}, title={#3}, title style={font=\small,
      yshift=8pt}, clip=false, legend style={font=\scriptsize, at={(0.5,-0.42)}, anchor=north, legend columns=4,
      draw=none, /tikz/every even column/.append style={column sep=6pt}}, legend cell align=left,
      ytick style={draw=none}, xtick style={draw=none}, ymajorgrids, grid style={cGray!25}]
      \addplot[draw=none, forget plot] table[x=pos, y expr=0] {#1};% the ticks (xtick=data)
      \namedbar{lvA}{lb12}{-0.325}\addlegendentry{LB $(1,2)$}
      \namedbar{lvB}{lb13}{-0.195}\addlegendentry{LB $(1,3)$}
      \namedbar{lvC}{lb22}{-0.065}\addlegendentry{LB $(2,2)$}
      \namedbar{lvD}{lb23}{0.065}\addlegendentry{LB $(2,3)$}
      \namedbar{lvE}{lb33}{0.195}\addlegendentry{LB $(3,3)$}
      \namedbar{lvF}{lb24}{0.325}\addlegendentry{LB $(2,4)$}
      \addplot[only marks, mark=diamond*, mark size=2.2pt, cB, mark options={draw=white, line width=0.4pt}]
        table[x expr=\thisrow{pos}-0.14, y expr=\thisrow{ub_extract}/\thisrow{ref}] {#1};
      \addlegendentry{UB: extracted ($k=m=3$)}
      \addplot[only marks, mark=x, mark size=2.4pt, black, thick]
        table[x expr=\thisrow{pos}+0.14, y expr=\thisrow{ub_hjw}/\thisrow{ref}] {#1};
      \addlegendentry{UB: HJW}
      \addplot[only marks, mark=none, nodes near coords={$\pgfplotspointmeta$}, point meta=explicit symbolic,
        nodes near coords style={anchor=south, font=\tiny, inner sep=1pt}, forget plot]
        table[x=pos, y=label_y, meta=ref_label] {#1};
      \draw[cGray] (axis cs:0.3,1) -- (axis cs:14.7,1);
      \draw[cGray, densely dashed] (axis cs:10,0) -- (axis cs:10,1.38);
      \node[font=\scriptsize, anchor=south] at (axis cs:5,1.38) {PPT (bound entangled)};
      \node[font=\scriptsize, anchor=south] at (axis cs:12.5,1.38) {NPT};
    \end{axis}
  \end{tikzpicture}}
%
% \polypanel{measure column prefix (EF or E1)}: the optimal polynomial of each state, and the points where it touches F
% (the Schmidt coefficients of the Wootters decomposition).
\providecommand{\polypanel}[1]{%
  \addplot[cA, thick] table[x=alpha, y=werner2] {figures/data/polynomials_#1.csv};\addlegendentry{Werner $2\times2$}
  \addplot[cA, only marks, mark=*, mark size=1.8pt, forget plot]
    table[x=alpha_werner2, y=#1_werner2] {figures/data/polynomials_touching.csv};
  \addplot[cB, thick] table[x=alpha, y=gisin] {figures/data/polynomials_#1.csv};\addlegendentry{Gisin}
  \addplot[cB, only marks, mark=*, mark size=1.8pt, forget plot]
    table[x=alpha_gisin, y=#1_gisin] {figures/data/polynomials_touching.csv};
  \addplot[cC, thick] table[x=alpha, y=mems] {figures/data/polynomials_#1.csv};\addlegendentry{MEMS}
  \addplot[cC, only marks, mark=*, mark size=1.8pt, forget plot]
    table[x=alpha_mems, y=#1_mems] {figures/data/polynomials_touching.csv};
  \addplot[cD, thick] table[x=alpha, y=bellprod] {figures/data/polynomials_#1.csv};\addlegendentry{Bell + $\ket{01}$}
  \addplot[cD, only marks, mark=*, mark size=1.8pt, forget plot]
    table[x=alpha_bellprod, y=#1_bellprod] {figures/data/polynomials_touching.csv};}%
\begin{figure*}
  \centering
  \begin{tikzpicture}
    \begin{groupplot}[group style={group size=2 by 1, horizontal sep=1.6cm}, width=0.46\textwidth,
      height=0.32\textwidth, tick label style={font=\footnotesize}, label style={font=\small}, title style={font=\small},
      xlabel={$\alpha_1$}, ylabel={$F(\alpha)$, $P(\alpha)$}, xmin=0, xmax=1, table/col sep=comma,
      legend style={font=\scriptsize, at={(0.5,0.02)}, anchor=south, draw=none, fill=white, fill opacity=0.85,
      text opacity=1}, legend cell align=left]
    \nextgroupplot[title={(a) Entanglement of formation, level $(3,4)$}, ymin=-0.8, ymax=1.15]
      \addplot[cGray, dashed, forget plot] coordinates {(0,0) (1,0)};
      \addplot[black, very thick] table[x=alpha, y=F] {figures/data/polynomials_EF.csv};\addlegendentry{$S(\alpha)$}
      \polypanel{EF}
    \nextgroupplot[title={(b) Vidal monotone $E_1$, level $(3,3)$}, ymin=-0.4, ymax=0.62]
      \addplot[cGray, dashed, forget plot] coordinates {(0,0) (1,0)};
      \addplot[black, very thick] table[x=alpha, y=F] {figures/data/polynomials_E1.csv};\addlegendentry{$E_1(\alpha)$}
      \polypanel{E1}
    \end{groupplot}
  \end{tikzpicture}
  \caption{
  Polynomials $P\leq F$ extracted from the SDP with embedded LP for different $2$-qubit states. For each state, this is the polynomial that yields the best bound on the convex roof on the grid, at level $(3,4)$ for $E_F$ (a) and at level $(3,3)$ for $E_1$ (b). Here we show
  that the LP is adaptive depending on the state and there is not a single polynomial that yields the best bound on the convex roof for all states at the same level. The states used are $\{$ Werner $2\times2$: singlet weight $p=0.8$; Gisin: $\lambda=0.9$, $\theta=\pi/8$; MEMS: $C=0.8$; Bell + $\ket{01}$: $p=0.6$ $\}$. The polynomials are shown as a function of the entanglement spectrum $\a = (\a_1, \a_2) = (\a_1, 1-\a_1)$.}
  \label{fig:polynomials}
\end{figure*}
